% GENERATED FILE. Edit canonical lessons, then run npm run generate:books.
% canonical-source-hash: fnv1a64:46d5f57bb839442a
% canonical-lessons: LA-C145-acerbus, LA-C145-vivus, LA-C145-mitis, LA-C145-ferox, LA-C145-iucundus

\chapter{Sour, Alive, Gentle, Fierce, Pleasant}
\label{ch:la-sour-alive-gentle-fierce-pleasant}
\hlchaptermodality{145}

\begin{chapteropening}
\textbf{By the end of this chapter:} \emph{I can say sour, alive, gentle, fierce and pleasant.}

Complete the last lesson of chapter 145: I can say sour, alive, gentle, fierce and pleasant.
\end{chapteropening}

\section[acerbus, acerba, acerbum]{acerbus, acerba, acerbum — sour, unripe}

\label{lesson:LA-C145-acerbus}



\begin{quote}
Before the new one: say the Latin for blind, then the Latin for sharp.
\end{quote}

\subsection*{You'll want to know: acerbus, acerba, acerbum}
\textbf{acerbus, acerba, acerbum} — \textquotedblleft{}sour, unripe\textquotedblright{}.

\begin{cousinweb}
\textbf{acerbus} is built on \emph{ācer}, \textquotedblleft{}sharp\textquotedblright{} (Lewis \& Short), like \emph{superbus}, \textquotedblleft{}proud\textquotedblright{}, on \emph{super}, \textquotedblleft{}above\textquotedblright{}. \emph{ācer} is usually traced to Proto-Indo-European \emph{*h2kros}, \textquotedblleft{}sharp\textquotedblright{} (EtymDB 2.1).
\end{cousinweb}

\begin{grammarlens}[title={What you've built}]
One more word for this chapter.
\end{grammarlens}

\subsection*{Your turn}
\begin{itemize}
\raggedright
  \item \emph{Say it:} \emph{acerbus}
  \item \emph{Say it:} \emph{acerbus}, once more
  \item \emph{Say it:} say \emph{acerbus}
  \item \emph{Recall:} say the Latin for blind, then the Latin for sharp, then say \emph{acerbus} again
\end{itemize}

\begin{tcolorbox}[breakable,colback=teal!4,colframe=teal!35!black,title={Before you move on}]
What does acerbus, acerba, acerbum mean? (\textquotedblleft{}Sour, unripe\textquotedblright{}.) Say it once more.
\end{tcolorbox}
\section[vīvus, vīva, vīvum]{vīvus, vīva, vīvum — alive}

\label{lesson:LA-C145-vivus}



\begin{quote}
Before the new one: say the Latin for sharp, then the Latin for sour.
\end{quote}

\subsection*{You'll want to know: vīvus, vīva, vīvum}
\textbf{vīvus, vīva, vīvum} — \textquotedblleft{}alive\textquotedblright{}.

\begin{cousinweb}
\textbf{vīvus} is inherited, usually traced to Proto-Indo-European \emph{*gwih3wos}, \textquotedblleft{}alive\textquotedblright{} (EtymDB 2.1), and belongs with \emph{vīvō}, \textquotedblleft{}to live\textquotedblright{} (Lewis \& Short). English \emph{quick}, once \textquotedblleft{}alive\textquotedblright{}, is a cousin.
\end{cousinweb}

\begin{grammarlens}[title={What you've built}]
One more word for this chapter.
\end{grammarlens}

\subsection*{Your turn}
\begin{itemize}
\raggedright
  \item \emph{Say it:} \emph{vīvus}
  \item \emph{Say it:} \emph{vīvus}, once more
  \item \emph{Say it:} say \emph{vīvus}
  \item \emph{Recall:} say the Latin for sharp, then the Latin for sour, then say \emph{vīvus} again
\end{itemize}

\begin{tcolorbox}[breakable,colback=teal!4,colframe=teal!35!black,title={Before you move on}]
What does vīvus, vīva, vīvum mean? (\textquotedblleft{}Alive\textquotedblright{}.) Say it once more.
\end{tcolorbox}
\section[mītis, mīte]{mītis, mīte — gentle, mild}

\label{lesson:LA-C145-mitis}



\begin{quote}
Before the new one: say the Latin for sour, then the Latin for alive.
\end{quote}

\subsection*{You'll want to know: mītis, mīte}
\textbf{mītis, mīte} — \textquotedblleft{}gentle, mild\textquotedblright{}.

\begin{cousinweb}
\textbf{mītis} is probably inherited, from a Proto-Indo-European \emph{*meh1y-}, \textquotedblleft{}mild, soft\textquotedblright{} (EtymDB 2.1); Lewis \& Short call it doubtful.
\end{cousinweb}

\begin{grammarlens}[title={What you've built}]
One more word for this chapter.
\end{grammarlens}

\subsection*{Your turn}
\begin{itemize}
\raggedright
  \item \emph{Say it:} \emph{mītis}
  \item \emph{Say it:} \emph{mītis}, once more
  \item \emph{Say it:} say \emph{mītis}
  \item \emph{Recall:} say the Latin for sour, then the Latin for alive, then say \emph{mītis} again
\end{itemize}

\begin{tcolorbox}[breakable,colback=teal!4,colframe=teal!35!black,title={Before you move on}]
What does mītis, mīte mean? (\textquotedblleft{}Gentle, mild\textquotedblright{}.) Say it once more.
\end{tcolorbox}
\section[ferōx, ferōcis]{ferōx, ferōcis — fierce}

\label{lesson:LA-C145-ferox}



\begin{quote}
Before the new one: say the Latin for alive, then the Latin for gentle.
\end{quote}

\subsection*{You'll want to know: ferōx, ferōcis}
\textbf{ferōx, ferōcis} — \textquotedblleft{}fierce\textquotedblright{}.

\begin{cousinweb}
\textbf{ferōx} is built on \emph{ferus}, \textquotedblleft{}wild\textquotedblright{} (Lewis \& Short), with \emph{-ōx}, an ending of strong character: wild-looking, fierce. \emph{ferus} is a cousin of Greek \emph{thēr}, \textquotedblleft{}wild beast\textquotedblright{} (Lewis \& Short).
\end{cousinweb}

\begin{grammarlens}[title={What you've built}]
One more word for this chapter.
\end{grammarlens}

\subsection*{Your turn}
\begin{itemize}
\raggedright
  \item \emph{Say it:} \emph{ferōx}
  \item \emph{Say it:} \emph{ferōx}, once more
  \item \emph{Say it:} say \emph{ferōx}
  \item \emph{Recall:} say the Latin for alive, then the Latin for gentle, then say \emph{ferōx} again
\end{itemize}

\begin{tcolorbox}[breakable,colback=teal!4,colframe=teal!35!black,title={Before you move on}]
What does ferōx, ferōcis mean? (\textquotedblleft{}Fierce\textquotedblright{}.) Say it once more.
\end{tcolorbox}
\section[iūcundus, iūcunda, iūcundum]{iūcundus, iūcunda, iūcundum — pleasant}

\label{lesson:LA-C145-iucundus}



\begin{quote}
Before the new one: say the Latin for gentle, then the Latin for fierce.
\end{quote}

\subsection*{You'll want to know: iūcundus, iūcunda, iūcundum}
\textbf{iūcundus, iūcunda, iūcundum} — \textquotedblleft{}pleasant\textquotedblright{}.

\begin{cousinweb}
\textbf{iūcundus} is usually derived from \emph{iuvō}, \textquotedblleft{}to help, to delight\textquotedblright{}; the link with \emph{iocus}, \textquotedblleft{}a joke\textquotedblright{}, that Lewis \& Short give is a later popular one.
\end{cousinweb}

\begin{grammarlens}[title={What you've built}]
One more word for this chapter.
\end{grammarlens}

\subsection*{Your turn}
\begin{itemize}
\raggedright
  \item \emph{Say it:} \emph{iūcundus}
  \item \emph{Say it:} \emph{iūcundus}, once more
  \item \emph{Say it:} say \emph{iūcundus}
  \item \emph{Recall:} say the Latin for gentle, then the Latin for fierce, then say \emph{iūcundus} again
\end{itemize}

\begin{tcolorbox}[breakable,colback=teal!4,colframe=teal!35!black,title={Before you move on}]
What does iūcundus, iūcunda, iūcundum mean? (\textquotedblleft{}Pleasant\textquotedblright{}.) Say it once more.
\end{tcolorbox}
